\section{Results}

\subsection{Image Census and Exclusions}

A total of 325 images at $t = 0$ images were processed across two exposure conditions
(Candesartan, $n$~=~191; ALK5i, $n$~=~134) and two independent experiments
(EXP1, $n$~=~158; EXP2, $n$~=~167).

Manual review identified \textbf{40 images (12.3\%)} containing no visible wound,
which were excluded from accuracy analysis.
The no-wound images were unevenly distributed (Table~\ref{tab:nowound-distribution}):
38 of 40 occurred in EXP1.
Two images remained unannotated.
The remaining \textbf{285 wound-bearing images} formed the validation cohort.

\begin{table}[htbp]
\centering
\caption{Distribution of no-wound images by condition and experiment.
  These images contained no visible scratch wound and were excluded from accuracy analysis.}
\label{tab:nowound-distribution}
\begin{tabular}{lrrr}
\toprule
Condition & EXP1 & EXP2 & Total \\
\midrule
Candesartan & 29 & 1 & 30 \\
ALK5i       &  9 & 1 & 10 \\
\midrule
\textbf{Total} & \textbf{38} & \textbf{2} & \textbf{40} \\
\bottomrule
\end{tabular}
\end{table}

\subsection{Wound Detection Accuracy}

Of 285 wound-bearing images, the algorithm correctly segmented \textbf{282 (98.9\%)} as judged
by expert review.
Only 3 images were incorrectly segmented.
Performance was consistent across conditions and experiments (Table~\ref{tab:accuracy}).

\begin{table}[htbp]
\centering
\caption{Detection accuracy by exposure and experiment (wound-bearing images only).
  Values are expressed as correctly segmented / total (percentage).
  Expert review served as ground truth.}
\label{tab:accuracy}
\begin{tabular}{llll}
\toprule
Condition & EXP1 & EXP2 & Total \\
\midrule
Candesartan & 67/67 (100.0\%) & 93/94 (98.9\%) & 160/161 (99.4\%) \\
ALK5i       & 51/53 (96.2\%)  & 71/71 (100.0\%) & 122/124 (98.4\%) \\
\midrule
\textbf{Total} & 118/120 (98.3\%) & 164/165 (99.4\%) & \textbf{282/285 (98.9\%)} \\
\bottomrule
\end{tabular}
\end{table}

\subsection{Automatic Quality-Control Evaluation}

The automatic QC flag was evaluated against human verdicts in two contexts.

\textbf{All annotated images ($n$~=~323).}
The QC algorithm achieved 94.4\% agreement with the human reviewer (Cohen's $\kappa$~=~0.69;
Table~\ref{tab:qc-confusion}).
It identified all 282 correctly segmented images (sensitivity~=~1.00, NPV~=~1.00) but
flagged only 23 of 41 problematic images as invalid (specificity~=~0.56).
The 18 false positives (QC valid, human incorrect) were predominantly no-wound images
(16 of 18) that the variance-based QC could not distinguish from valid segmentations.

\textbf{Wound-bearing images only ($n$~=~285).}
The QC algorithm correctly classified 283 of 285 images (accuracy~=~99.3\%),
missing 2 incorrect segmentations while producing no false negatives.

\begin{table}[htbp]
\centering
\caption{QC confusion matrix (all annotated images, $n$~=~323).
  Sensitivity~=~1.00; Specificity~=~0.56; PPV~=~0.94; NPV~=~1.00.}
\label{tab:qc-confusion}
\begin{tabular}{lrr}
\toprule
 & Human: Correct & Human: Incorrect \\
\midrule
\textbf{QC: Valid}   & 282 (TP) & 18 (FP) \\
\textbf{QC: Invalid} & 0 (FN)   & 23 (TN) \\
\bottomrule
\end{tabular}
\end{table}

\subsection{Pixel-Level Segmentation Quality}

For 5 images with expert-drawn polygon annotations, pixel-level metrics were computed against
human ground truth (Table~\ref{tab:pixel-metrics}).
The mean Dice coefficient was 0.91 (SD~=~0.06), indicating high spatial overlap.
The high recall (0.97) relative to precision (0.87) indicates a tendency to slightly
over-segment (include background pixels) rather than miss wound area.

\begin{table}[htbp]
\centering
\caption{Pixel-level segmentation metrics against expert polygon annotations ($n$~=~5).
  Dice and F1 are identical for binary masks.
  Hausdorff distance is reported at the 95th percentile.}
\label{tab:pixel-metrics}
\begin{tabular}{l
  S[table-format=1.3]
  S[table-format=1.3]
  S[table-format=1.3]
  S[table-format=1.3]}
\toprule
{Metric} & {Mean} & {SD} & {Min} & {Max} \\
\midrule
Dice                   & 0.909 & 0.057 & 0.843 & 0.998 \\
IoU                    & 0.837 & 0.099 & 0.729 & 0.997 \\
Precision              & 0.868 & 0.123 & 0.729 & 1.000 \\
Recall                 & 0.969 & 0.070 & 0.844 & 1.000 \\
F1                     & 0.909 & 0.057 & 0.843 & 0.998 \\
Accuracy               & 0.954 & 0.029 & 0.921 & 0.999 \\
{Hausdorff 95 (px)}    & 45.6  & 27.8  & 0.0   & 76.0  \\
\bottomrule
\end{tabular}
\end{table}

\subsection{Wound Healing Dynamics}

\subsubsection{Validated Cohort and Pipeline Illustration}

Of the 282 verified segmentations, all were included in wound healing dynamics analysis:
DMSO ($n$~=~100), Candesartan ($n$~=~86), ALK5i ($n$~=~55), Media ($n$~=~41).
Two cell lines were represented: iMC~ISOR544C ($n$~=~145) and iMC~MUTR544C ($n$~=~137).
Each trajectory comprised 20--25 timepoints over approximately 24~hours.
A representative pipeline output is shown in Figure~\ref{fig:fig1}.

\begin{figure}[htbp]
  \centering
  \includegraphics[width=\linewidth]{fig1_model_quality}
  \caption{%
    \textbf{Automated wound boundary detection and model quality metrics.}
    \textbf{(A--C)} Representative phase-contrast images of a wound healing assay at three
    timepoints ($t = 0$, mid, and $t = 20$~h) with computationally detected wound boundaries
    overlaid as semi-transparent filled regions.
    Earlier timepoints are shown as faint dashed outlines; the current timepoint as solid fill.
    Images were acquired at 1.24~\textmu m/px.
    \textbf{(D)} Normalised wound area over time, expressed as a fraction of the $t = 0$
    wound area.
    Error bands represent 95\% confidence intervals.
    Sample sizes per group are provided in the legend.
    (Figure~\ref{fig:fig1} corresponds to the analysis of the ALK5i exposure,
    one cell line subset; full dataset is shown in Figure~\ref{fig:fig2}.)%
  }
  \label{fig:fig1}
\end{figure}

\subsubsection{Fractional Closure at Fixed Timepoints}

Mean fractional closure at 24~hours ranged from 15.5\% (ALK5i) to 20.0\%
(Candesartan), with DMSO at 19.7\% and Media at 18.5\% (Table~\ref{tab:closure24h}).
At 12~hours, the pattern was similar: ALK5i showed the lowest mean closure (12.0\%) and
DMSO/Candesartan were comparable (14.5\% and 14.9\%; Table~\ref{tab:closure12h}).

\begin{table}[htbp]
\centering
\caption{Mean fractional wound closure at 24~hours by treatment condition (mean~$\pm$~SEM,
  all concentrations pooled).
  No condition reached 50\% closure within the observation window.}
\label{tab:closure24h}
\begin{tabular}{l
  r
  S[table-format=1.3]
  S[table-format=1.3]}
\toprule
{Condition} & {$n$} & {Mean Closure} & {SEM} \\
\midrule
DMSO        & 100 & 0.197 & 0.010 \\
Candesartan &  86 & 0.200 & 0.011 \\
ALK5i       &  55 & 0.155 & 0.008 \\
Media       &  41 & 0.185 & 0.012 \\
\bottomrule
\end{tabular}
\end{table}

\begin{table}[htbp]
\centering
\caption{Mean fractional wound closure at 12~hours by treatment condition (mean~$\pm$~SEM,
  all concentrations pooled).}
\label{tab:closure12h}
\begin{tabular}{l
  r
  S[table-format=1.3]
  S[table-format=1.3]}
\toprule
{Condition} & {$n$} & {Mean Closure} & {SEM} \\
\midrule
DMSO        & 100 & 0.145 & 0.006 \\
Candesartan &  86 & 0.149 & 0.009 \\
ALK5i       &  55 & 0.120 & 0.007 \\
Media       &  41 & 0.131 & 0.008 \\
\bottomrule
\end{tabular}
\end{table}

\subsubsection{Time to 50\% Closure}

No condition reached 50\% closure within the 24-hour observation window.

\subsubsection{Time-Course Profiles and Cross-Sectional Analysis}

Wound closure trajectories followed a characteristic pattern across all conditions
(Figure~\ref{fig:fig2}):
an initial lag phase (0--4~h) with minimal closure, followed by a progressive closure phase
that did not plateau within 24~hours.
The ALK5i-treated group showed the slowest kinetics throughout the observation period,
separating from other conditions by approximately 6--8~hours post-wounding.
DMSO and candesartan groups showed similar kinetic profiles, while the media control fell
between ALK5i and the vehicle-treated groups.

\begin{figure}[htbp]
  \centering
  \includegraphics[width=\linewidth]{fig2_wound_dynamics}
  \caption{%
    \textbf{Wound healing dynamics across experimental conditions.}
    \textbf{(A)} Mean wound closure speed (pixels/h) at $t = 20$~h,
    computed as the rate of wound edge displacement.
    \textbf{(B)} Mean wound closure distance (pixels) at $t = 20$~h,
    representing cumulative inward displacement of the wound boundary from $t = 0$.
    \textbf{(C)} Time-series of wound edge displacement over the full observation period.
    Error bars represent 95\% confidence intervals.
    Statistical comparisons were performed using Welch's $t$-test with Bonferroni correction.
    Sample sizes: Line~1~--~Alk5i: $n = 10$; Line~1~--~DMSO: $n = 8$;
    Line~1~--~Media: $n = 15$; Line~2~--~Alk5i: $n = 5$;
    Line~2~--~DMSO: $n = 15$; Line~2~--~Media: $n = 6$.%
  }
  \label{fig:fig2}
\end{figure}

Sample-level accounting, QC pass rates, and full descriptive statistics are provided in
Supplementary Tables~\ref{tab:supp-sample-sizes},
\ref{tab:supp-sample-accounting}, and~\ref{tab:supp-qc-validation}.
Inter-experiment reproducibility data are reported in Supplementary
Table~\ref{tab:supp-reproducibility}.
Full pairwise comparison results including confidence intervals appear in
Supplementary Table~\ref{tab:supp-pairwise}.
The full descriptive statistics per group are in Supplementary
Table~\ref{tab:supp-descriptive}.

Cross-sectional analysis at the final timepoint confirmed no statistically significant
differences between ALK5i and DMSO after Bonferroni correction across both cell lines
(Line~1: $p = 0.38$, Cohen's $d = -0.40$; Line~2: $p = 0.40$, Cohen's $d = -0.29$;
Table~\ref{tab:effects}).
Media did not differ significantly from DMSO in either cell line ($p > 0.59$ in both cases).
These results are consistent with modest absolute differences that are underpowered given
current sample sizes.

\begin{table}[htbp]
\centering
\caption{Effect sizes and pairwise comparisons vs.\ DMSO at the final timepoint.
  Welch's $t$-test with Bonferroni correction; $p$-values reported after correction.
  All comparisons non-significant after correction.}
\label{tab:effects}
\begin{tabular}{ll
  S[table-format=2.2]
  S[table-format=2.2]
  S[table-format=2.2]
  S[table-format=-1.3]
  S[table-format=1.3]
  rr}
\toprule
{Cell line} & {Treatment} & {Control mean} & {Treated mean} & {Difference} & {Cohen's $d$} & {$p$ (corrected)} & {$n_\text{ctrl}$} & {$n_\text{trt}$} \\
\midrule
Line 1 & Alk5i & 19.28 & 16.34 & -2.94 & -0.402 & 1.000 & 8 & 10 \\
Line 1 & Media & 19.28 & 20.87 &  1.59 &  0.200 & 1.000 & 8 & 15 \\
Line 2 & Alk5i & 21.69 & 16.56 & -5.13 & -0.294 & 1.000 & 15 &  5 \\
Line 2 & Media & 21.69 & 18.86 & -2.83 & -0.169 & 1.000 & 15 &  6 \\
\bottomrule
\end{tabular}
\smallskip

\noindent\textit{Note:} Mean distances are in pixels.
All $p$-values~=~1.000 after Bonferroni correction (4 comparisons).
Cohen's $d$ is computed using pooled standard deviation.
\end{table}

\subsubsection{Mixed-Effects Time-Series Model}

The linear mixed-effects model (Equation~\eqref{eq:mixed-model}) confirmed a highly
significant main effect of time ($\beta = 0.879$~px/frame, $p < 10^{-12}$),
indicating consistent wound closure across all conditions.
No time $\times$ condition or time $\times$ cell line interaction reached significance
(all $p > 0.31$; Supplementary Table~\ref{tab:supp-mixed-model-coefficients}),
consistent with the cross-sectional analysis.
Full mixed-model parameter estimates and variance components are provided in
Supplementary Table~\ref{tab:supp-mixed-model-coefficients}.
